% !TeX root = main.tex

\phantomsection
\chapter*{译者说明}
\addcontentsline{toc}{chapter}{译者说明}

本书译自 Rosalie Chevalley 的维也纳大学理学院学期项目论文(Semester Project, Winter 2009--2010)
\emph{Profinite groups and Galois cohomology}，指导教师为 Dietrich Burde，
助教为 Wolfgang Moens. 翻译与 LaTeX 重排以原书 PDF 为最高依据，
MinerU 与 PaddleOCR 的提取文本仅用作工作底稿.

排版遵循以下对应关系：

\begin{enumerate}
    \item 章、节、小节的题号与原题一致：节号不带章前缀（如``1. 定义与例子''、``2.3 性质''）；
    \item 定理类环境按节编号，且与原书一样分为两组各自共用计数器：
          定义、记号、注、例共用一个计数器，定理、引理、命题、推论共用另一个计数器.
          因此同一节中可能出现``命题 2.4 与定义 2.3 相邻''的情形，这是原书的编号方式，并非排印错误；
    \item 行间公式沿用原书全文连续的编号 (1)--(7)；
    \item 原书正文中方括号式引文（如 [Rib99]）全部改为可点击的交叉引用，
          文献细节（定理号、章节号）保留在引文注记中.
\end{enumerate}

原书为学位论文，个别处编号与语法存在笔误，译本一律照原样保留，不作擅自改动；
涉及数学表述的明显讹误则在脚注中说明. 数学术语的译法参照 Ribes 原著中译本
与国内同调代数、域论论著的习惯译名，首次出现时以楷体标出并附英文原词，
全书收录于文末名词索引.

\vspace{0.5cm}
\hfill \Translator

\hfill \TranslationDate
